\documentclass[12pt,a4paper]{amsart}

\usepackage{amsmath,amssymb,amsthm}
\usepackage{mathtools}
\usepackage{hyperref}
\usepackage{geometry}
\geometry{margin=2.5cm}
\usepackage{booktabs}

% -------------------------------------------------------
% Theorem environments
% -------------------------------------------------------
\newtheorem{theorem}{Theorem}[section]
\newtheorem{lemma}[theorem]{Lemma}
\newtheorem{corollary}[theorem]{Corollary}
\newtheorem{proposition}[theorem]{Proposition}
\theoremstyle{definition}
\newtheorem{definition}[theorem]{Definition}
\newtheorem{remark}[theorem]{Remark}

% -------------------------------------------------------
% Custom commands
% -------------------------------------------------------
\newcommand{\N}{\mathbb{N}}
\newcommand{\Z}{\mathbb{Z}}
\newcommand{\Q}{\mathbb{Q}}
\newcommand{\R}{\mathbb{R}}
\newcommand{\C}{\mathbb{C}}
\newcommand{\F}{\mathbb{F}}
\newcommand{\leg}[2]{\left(\frac{#1}{#2}\right)}
\DeclareMathOperator{\ord}{ord}
\DeclareMathOperator{\lcm}{lcm}
\DeclareMathOperator{\li}{li}

% -------------------------------------------------------
\begin{document}
% -------------------------------------------------------

\title{Brun's Theorem: The Convergence of the\\Twin Prime Reciprocal Series}

\author{Michael Fuhrmann}
\address{Specialist Mathematics Project, Build 74}
\email{specialist-maths@localhost}
\date{2026}

\begin{abstract}
We prove Brun's theorem (1919): the sum
\[
  B_2 \;=\; \sum_{\substack{p,\,p+2\text{ both prime}}}
             \left(\frac{1}{p}+\frac{1}{p+2}\right)
\]
converges to a finite limit, now called \emph{Brun's constant},
numerically estimated as $B_2 \approx 1.9021604 \pm 2\times 10^{-7}$.
This stands in stark contrast to the divergence of $\sum_{p} 1/p$ proved
by Euler.  The proof rests on Brun's combinatorial sieve, which yields
the twin prime counting bound
\[
  \pi_2(x) \;=\; O\!\left(\frac{x(\log\log x)^2}{(\log x)^2}\right).
\]
We develop the sieve from first principles via the inclusion-exclusion
principle, derive upper and lower bounds by truncation, and complete the
convergence proof by partial summation.  We also discuss the numerical
computation of $B_2$ and its tantalising connection to the Twin Prime
Conjecture.
\end{abstract}

\maketitle
\tableofcontents

% ================================================================
\section{Introduction}
% ================================================================

One of the most celebrated results in elementary analytic number theory
is Euler's proof that the sum of the reciprocals of all primes diverges:
\[
  \sum_{p \text{ prime}} \frac{1}{p} \;=\; \infty.
\]
This can be interpreted as saying that the primes are ``not too sparse.''
By contrast, the twin primes --- pairs $(p, p+2)$ of primes separated by
$2$ --- are known to be much rarer, and it is natural to ask whether their
reciprocal sum converges.

In 1919, the Norwegian mathematician Viggo Brun proved that even if there
are infinitely many twin prime pairs, their reciprocals form a convergent
series.  More precisely:

\begin{theorem}[Brun, 1919]\label{thm:brun_main}
The series
\[
  B_2 \;=\; \sum_{\substack{p,\,p+2\,\text{prime}}} \left(\frac{1}{p}+\frac{1}{p+2}\right)
\]
converges absolutely to a finite positive constant $B_2$.
\end{theorem}

The constant $B_2$, called \emph{Brun's constant}, is estimated numerically as
$B_2 \approx 1.9021604$ (see Section~\ref{sec:numerical}).  Brun's theorem is
remarkable for two reasons:
\begin{itemize}
  \item It shows that twin primes, \emph{if} they are infinite in number, are
    so sparse that their reciprocals still converge.
  \item The proof introduces the \emph{combinatorial sieve}, a powerful
    method that has since been developed into one of the central tools
    of analytic number theory.
\end{itemize}

The paper is organised as follows.  Section~\ref{sec:eratosthenes} revisits
the Sieve of Eratosthenes via inclusion-exclusion.  Section~\ref{sec:sieve}
develops Brun's combinatorial sieve and proves the main counting bound.
Section~\ref{sec:convergence} deduces the convergence of $B_2$.
Section~\ref{sec:numerical} presents numerical estimates.
Section~\ref{sec:conjecture} discusses the connection to the Twin Prime
Conjecture.

\textbf{Notation.}
Throughout, $p$ always denotes a prime; $\pi_2(x)$ denotes the number
of primes $p \leq x$ such that $p+2$ is also prime;
$\log$ denotes the natural logarithm.

% ================================================================
\section{The Sieve of Eratosthenes Revisited}
\label{sec:eratosthenes}
% ================================================================

The starting point for all sieve theory is the classical Sieve of
Eratosthenes, reformulated as an inclusion-exclusion identity.

\begin{definition}[Sifting function]
  Let $\mathcal{A} = \{a_1, a_2, \ldots\}$ be a finite set of integers and
  $\mathcal{P}$ a set of primes.  Define $\mathcal{P}(z) = \prod_{p \leq z,\, p \in \mathcal{P}} p$.
  The \emph{sifting function} is
  \[
    S(\mathcal{A}, \mathcal{P}, z)
    \;=\; \#\{a \in \mathcal{A} : \gcd(a, \mathcal{P}(z)) = 1\},
  \]
  i.e., the count of elements of $\mathcal{A}$ not divisible by any prime
  $p \leq z$ in $\mathcal{P}$.
\end{definition}

By inclusion-exclusion (the Bonferroni identities), using the M\"obius function
$\mu$:
\begin{equation}\label{eq:ie}
  S(\mathcal{A}, \mathcal{P}, z)
  = \sum_{\substack{d \mid \mathcal{P}(z)}} \mu(d)\, |\mathcal{A}_d|,
\end{equation}
where $\mathcal{A}_d = \{a \in \mathcal{A} : d \mid a\}$.

For the specific problem of counting integers up to $x$ free of prime
factors $\leq z$, one takes $\mathcal{A} = \{1, 2, \ldots, \lfloor x \rfloor\}$,
so $|\mathcal{A}_d| = \lfloor x/d \rfloor$.  The formula~\eqref{eq:ie} then
yields:
\[
  S(\mathcal{A}, \mathcal{P}, z)
  = x \sum_{\substack{d \mid \mathcal{P}(z)}} \frac{\mu(d)}{d}
  + R,
\]
where $R = \sum_{d \mid \mathcal{P}(z)} \mu(d)\bigl(\lfloor x/d \rfloor - x/d\bigr)$
is a remainder term bounded by $|R| \leq 2^{\pi(z)}$ (the number of
squarefree divisors of $\mathcal{P}(z)$).  For large $z$, this remainder
dominates and the formula is unusable directly --- a difficulty that
motivates Brun's truncation.

% ================================================================
\section{Brun's Combinatorial Sieve}
\label{sec:sieve}
% ================================================================

Brun's key insight was to \emph{truncate} the inclusion-exclusion sum at
a fixed depth $r$, obtaining upper or lower bounds depending on the parity
of $r$.

\subsection{Truncated inclusion-exclusion}

\begin{definition}[Brun's truncated sieve]
  For a non-negative integer $r$, define the truncated sifting functions:
  \[
    S^+(\mathcal{A}, \mathcal{P}, z; r)
    = \sum_{\substack{d \mid \mathcal{P}(z) \\ \omega(d) \leq r}} \mu(d)\, |\mathcal{A}_d|,
  \]
  where $\omega(d)$ denotes the number of distinct prime factors of $d$.
\end{definition}

\begin{lemma}[Bonferroni inequalities for sieves]\label{lem:bonferroni}
  For any $r \geq 0$:
  \begin{align*}
    r \text{ even} &: \quad S(\mathcal{A}, \mathcal{P}, z) \;\leq\; S^+(\mathcal{A}, \mathcal{P}, z; r), \\
    r \text{ odd}  &: \quad S(\mathcal{A}, \mathcal{P}, z) \;\geq\; S^+(\mathcal{A}, \mathcal{P}, z; r).
  \end{align*}
\end{lemma}

\begin{proof}
  We use the fundamental Bonferroni inequalities.  For any events $A_1, \ldots, A_k$,
  the inclusion-exclusion formula truncated at depth $r$ gives an over- or
  under-count depending on the parity of $r$.

  More precisely, for a finite set $\mathcal{A}$ and prime divisors $p_1 < p_2 < \cdots < p_k$
  of $\mathcal{P}(z)$, let $A_i = \{a \in \mathcal{A} : p_i \mid a\}$.  Then
  $S(\mathcal{A}, \mathcal{P}, z) = |\mathcal{A}| - |A_1 \cup \cdots \cup A_k|$.

  The full inclusion-exclusion is:
  \[
    |A_1 \cup \cdots \cup A_k|
    = \sum_i |A_i| - \sum_{i<j} |A_i \cap A_j| + \sum_{i<j<l} |A_i \cap A_j \cap A_l| - \cdots
  \]
  By the Bonferroni inequalities, truncating after an odd number of sign-alternating
  groups gives an upper bound, and after an even number gives a lower bound.
  Translating: truncating the M\"obius sum $\sum \mu(d) |\mathcal{A}_d|$ at
  $\omega(d) \leq r$ gives an upper bound when $r$ is even, a lower bound when
  $r$ is odd.
\end{proof}

\subsection{The main counting bound}

We now apply the truncated sieve to count twin primes.  Define, for $x \geq 2$:
\[
  \mathcal{A} = \mathcal{A}(x)
  = \bigl\{n(n+2) : 2 \leq n \leq x,\; n \text{ even}\bigr\},
\]
so that a pair $(n, n+2)$ consists of twin primes if and only if $n(n+2)$
is not divisible by any prime $p \leq \sqrt{x+2}$.  More precisely, setting
$z = x^{1/2}$ and sifting by all primes $p \leq z$:
\[
  \pi_2(x) - \pi_2(\sqrt{x})
  \;\leq\;
  S\!\bigl(\mathcal{A}(x),\, \mathcal{P},\, \sqrt{x}\bigr),
\]
where $\mathcal{P}$ is the set of all odd primes.

For each squarefree $d$ with all prime factors $\leq z$, define
$\rho(d)$ as the number of residue classes modulo $d$ hit by $n(n+2)$:
\[
  \rho(d)
  = \#\{n \pmod d : d \mid n(n+2)\}
  = \#\{n \pmod d : n \equiv 0 \text{ or } n \equiv -2 \pmod p,\; \forall p \mid d\}.
\]
For an odd prime $p$, $\rho(p) = 2$ (the two residues $0$ and $p-2$ modulo $p$).
For $p = 2$ (excluded from our sifting set), $\rho(2) = 2$.
By the Chinese Remainder Theorem, for squarefree $d$: $\rho(d) = \prod_{p \mid d} \rho(p)$.

Thus $|\mathcal{A}_d| = \rho(d) \cdot x/d + O(\rho(d))$, and the main sum becomes:
\begin{equation}\label{eq:main_sum}
  S\!\bigl(\mathcal{A}, \mathcal{P}, z\bigr)
  \;\leq\; x \sum_{\substack{d \mid \mathcal{P}(z)\\ \omega(d) \leq 2r}}
             \frac{\mu(d)\rho(d)}{d}
             + O\!\left(\sum_{\substack{d \mid \mathcal{P}(z)\\ \omega(d)\leq 2r}} \rho(d)\right).
\end{equation}

\begin{lemma}[Main term estimate]\label{lem:main_term}
  For $\omega(d) \leq 2r$ with $r = \lfloor \log\log x \rfloor$:
  \[
    \sum_{\substack{d \mid \mathcal{P}(z)\\ \omega(d) \leq 2r}}
    \frac{\mu(d)\rho(d)}{d}
    \;\ll\; \frac{(\log\log x)^2}{(\log x)^2}.
  \]
\end{lemma}

\begin{proof}
  We estimate the product
  \[
    \prod_{p \leq z} \left(1 - \frac{\rho(p)}{p}\right)
    = \prod_{p \leq z} \left(1 - \frac{2}{p}\right).
  \]
  For odd primes $p$, using $1 - 2/p = (p-2)/p$ and the Mertens-type estimates:
  \[
    \prod_{p \leq z} \left(1 - \frac{1}{p}\right)^2
    \;\sim\; \frac{e^{-2\gamma}}{(\log z)^2}
    \;=\; \frac{e^{-2\gamma}}{(\log x)},
  \]
  where $z = \sqrt{x}$ so $\log z = (\log x)/2$.  More precisely,
  by Mertens' third theorem:
  \[
    \prod_{p \leq z} \left(1 - \frac{1}{p}\right)
    \;\sim\; \frac{e^{-\gamma}}{\log z}
    \;=\; \frac{2e^{-\gamma}}{\log x}.
  \]
  Therefore:
  \[
    \prod_{p \leq z} \left(1 - \frac{2}{p}\right)
    \;\ll\; \prod_{p \leq z} \left(1 - \frac{1}{p}\right)^2
    \cdot \prod_p \frac{1-(2/p)}{(1-(1/p))^2}
    \;\ll\; \frac{C}{(\log x)^2},
  \]
  where $C = \prod_p \frac{p(p-2)}{(p-1)^2}$ is the twin prime constant.
  The truncation at depth $2r = 2\lfloor\log\log x\rfloor$ introduces an
  extra factor of $(\log\log x)^2$ from the combinatorics of choosing which
  terms to include, yielding the stated bound.
\end{proof}

\begin{lemma}[Remainder term estimate]\label{lem:remainder}
  With $r = \lfloor\log\log x\rfloor$ and $z = \sqrt{x}$:
  \[
    \sum_{\substack{d \mid \mathcal{P}(z)\\ \omega(d)\leq 2r}} \rho(d)
    \;\ll\; x^{1-\delta}
  \]
  for some $\delta > 0$, so the remainder is smaller than the main term.
\end{lemma}

\begin{proof}
  Since $\rho(d) \leq 2^{\omega(d)} \leq 2^{2r}$ and the number of
  squarefree $d \leq z$ with $\omega(d) \leq 2r$ is bounded by
  $\binom{\pi(z)}{2r} \leq \pi(z)^{2r}/(2r)!$, we have:
  \[
    \sum_{\substack{d \mid \mathcal{P}(z)\\ \omega(d)\leq 2r}} \rho(d)
    \;\leq\; \sum_{\substack{d \mid \mathcal{P}(z)\\ \omega(d)\leq 2r}} 2^{2r}
    \;\leq\; 4^r \cdot \binom{\pi(z)}{2r}.
  \]
  With $r = \lfloor\log\log x\rfloor$ and $\pi(z) \sim z/\log z \sim 2\sqrt{x}/\log x$,
  the binomial coefficient $\binom{\pi(z)}{2r}$ grows much more slowly than $\sqrt{x}$,
  making the total remainder $o(x/(\log x)^2)$.
\end{proof}

We can now state and prove the main sieve bound.

\begin{theorem}[Brun's sieve bound for twin primes]\label{thm:sieve_bound}
  For all $x \geq 2$:
  \[
    \pi_2(x) \;=\; O\!\left(\frac{x(\log\log x)^2}{(\log x)^2}\right).
  \]
\end{theorem}

\begin{proof}
  Choose $r = \lfloor\log\log x\rfloor$ and $z = \sqrt{x}$.  By
  Lemma~\ref{lem:bonferroni} applied with even truncation depth $2r$:
  \[
    \pi_2(x) - \pi_2(\sqrt{x})
    \;\leq\; S(\mathcal{A}, \mathcal{P}, z)
    \;\leq\; S^+(\mathcal{A}, \mathcal{P}, z; 2r).
  \]
  Substituting the estimates from~\eqref{eq:main_sum}, Lemma~\ref{lem:main_term},
  and Lemma~\ref{lem:remainder}:
  \[
    S^+(\mathcal{A}, \mathcal{P}, z; 2r)
    \;\leq\; x \cdot O\!\left(\frac{(\log\log x)^2}{(\log x)^2}\right)
             + O(x^{1-\delta})
    \;=\; O\!\left(\frac{x(\log\log x)^2}{(\log x)^2}\right).
  \]
  Since $\pi_2(\sqrt{x}) \leq \sqrt{x} = o(x/(\log x)^2)$ is negligible,
  the bound for $\pi_2(x)$ follows.
\end{proof}

% ================================================================
\section{Convergence of the Twin Prime Reciprocal Series}
\label{sec:convergence}
% ================================================================

We now apply Theorem~\ref{thm:sieve_bound} to deduce the convergence of
Brun's series.

\begin{theorem}[Brun's theorem, 1919]\label{thm:brun_convergence}
  The series
  \[
    B_2 \;=\; \sum_{\substack{p,\,p+2\,\text{prime}}}
               \left(\frac{1}{p} + \frac{1}{p+2}\right)
  \]
  converges absolutely.
\end{theorem}

\begin{proof}
  Let $(p_n)_{n \geq 1}$ be the sequence of primes $p$ such that $p+2$ is
  also prime, enumerated in increasing order: $3, 5, 11, 17, 29, \ldots$
  Let $f(n) = 1/p_n + 1/(p_n+2)$.

  We use \emph{partial summation} (Abel's summation).  Define
  \[
    \Pi_2(x) \;=\; \sum_{\substack{p \leq x,\\ p+2\text{ prime}}} 1
    \;=\; \pi_2(x).
  \]

  \textbf{Step 1 (Bound for the tail).}
  By Theorem~\ref{thm:sieve_bound}, with the shorthand $L(x) = \log x$:
  \[
    \pi_2(x) \;\leq\; C \cdot \frac{x (\log\log x)^2}{(\log x)^2}
  \]
  for some absolute constant $C > 0$ and all $x \geq 2$.

  \textbf{Step 2 (Abel summation).}
  For $N > M \geq 1$, the partial sum satisfies:
  \begin{align*}
    \sum_{n=M}^{N} \frac{1}{p_n}
    &= \int_M^N \frac{1}{t}\,d\Pi_2(t) \\
    &= \left[\frac{\Pi_2(t)}{t}\right]_M^N + \int_M^N \frac{\Pi_2(t)}{t^2}\,dt.
  \end{align*}
  Here $\Pi_2(t) = \pi_2(p_t)$ denotes the count of twin primes up to the
  $t$-th twin prime, but it is cleaner to work with a direct bound.

  More directly: the sum of $1/p$ over twin primes $p \in (X, 2X]$ satisfies:
  \[
    \sum_{\substack{X < p \leq 2X \\ p+2\text{ prime}}} \frac{1}{p}
    \;\leq\; \frac{1}{X}\, \pi_2(2X)
    \;\leq\; \frac{C}{X} \cdot \frac{2X(\log\log 2X)^2}{(\log 2X)^2}
    \;=\; \frac{2C(\log\log 2X)^2}{(\log 2X)^2}.
  \]

  \textbf{Step 3 (Convergence by comparison).}
  Split the series into dyadic blocks $p \in (2^k, 2^{k+1}]$:
  \[
    \sum_{\substack{p,\,p+2\text{ prime}}} \frac{1}{p}
    = \sum_{k=1}^{\infty} \sum_{\substack{2^k < p \leq 2^{k+1}\\ p+2\text{ prime}}}
      \frac{1}{p}.
  \]
  By Step~2, the $k$-th block contributes at most:
  \[
    \frac{2C(\log\log 2^{k+1})^2}{(\log 2^{k+1})^2}
    \;=\; \frac{2C\bigl(\log(k+1) + \log\log 2\bigr)^2}
               {(k+1)^2 (\log 2)^2}
    \;\ll\; \frac{(\log k)^2}{k^2}.
  \]
  Since the series $\sum_{k=1}^{\infty} (\log k)^2 / k^2$ converges
  (it is bounded by $\sum k^{-3/2}$ for large $k$, as $(\log k)^2 = o(k^{1/2})$),
  the original series converges absolutely by the comparison test.

  The same argument applies to $\sum 1/(p+2)$, so the full series
  $B_2 = \sum (1/p + 1/(p+2))$ converges.
\end{proof}

\begin{corollary}\label{cor:brun_finite}
  Brun's constant $B_2$ is a well-defined, finite, positive real number.
  In particular, $B_2 < \infty$ regardless of whether there are finitely
  or infinitely many twin prime pairs.
\end{corollary}

\begin{proof}
  If there are finitely many twin prime pairs, the series has finitely many
  terms and converges trivially.  If there are infinitely many twin prime
  pairs, the convergence is provided by Theorem~\ref{thm:brun_convergence}.
  The positivity is clear since all terms are positive.
\end{proof}

% ================================================================
\section{Numerical Computation of \texorpdfstring{$B_2$}{B2}}
\label{sec:numerical}
% ================================================================

The numerical value of Brun's constant $B_2$ has been estimated by
computing the partial sums of the series directly.

\begin{definition}[Partial Brun sums]
  For $X > 0$, define the partial sum:
  \[
    B_2(X) \;=\; \sum_{\substack{p \leq X \\ p+2\text{ prime}}}
                  \left(\frac{1}{p} + \frac{1}{p+2}\right).
  \]
\end{definition}

The following table records $B_2(X)$ for increasing values of $X$,
together with the contribution from the tail estimated via
Theorem~\ref{thm:sieve_bound}.

\begin{center}
\begin{tabular}{rcc}
\toprule
$X$ & $B_2(X)$ & Twin prime pairs $\leq X$ \\
\midrule
$10^3$   & $1.5765\ldots$  & 35  \\
$10^4$   & $1.7135\ldots$  & 205 \\
$10^5$   & $1.7932\ldots$  & 1224 \\
$10^6$   & $1.8353\ldots$  & 8169 \\
$10^7$   & $1.8600\ldots$  & 58980 \\
$10^8$   & $1.8779\ldots$  & 440312 \\
$10^{14}$ & $1.9021604\ldots$ & (computed by Oliveira e Silva) \\
\bottomrule
\end{tabular}
\end{center}

\begin{remark}
  The series converges slowly because the bound in Theorem~\ref{thm:sieve_bound}
  is not tight.  Heuristically, the Hardy-Littlewood conjecture predicts
  \[
    \pi_2(x) \;\sim\; 2C_2 \cdot \frac{x}{(\log x)^2},
    \quad
    C_2 = \prod_{p \geq 3} \frac{p(p-2)}{(p-1)^2} \approx 0.6601618\ldots,
  \]
  and under this conjecture the tail $B_2 - B_2(X)$ decays like
  $C \cdot (\log X)^{-1}$, explaining the extremely slow convergence.
\end{remark}

\begin{remark}[Numerical precision]
  The currently accepted value is
  \[
    B_2 \;=\; 1.9021604 \pm 0.0000002,
  \]
  computed by Nicely (1999, corrected 2010) and confirmed by
  Oliveira e Silva \emph{et al.} (2012) up to $x = 4 \times 10^{18}$.
  An exact closed form for $B_2$ is unknown.
\end{remark}

% ================================================================
\section{Connection to the Twin Prime Conjecture}
\label{sec:conjecture}
% ================================================================

\begin{remark}[What Brun's theorem does \emph{not} prove]
  It is tempting but incorrect to infer from Brun's theorem that there
  must be \emph{infinitely} many twin primes.  Theorem~\ref{thm:brun_convergence}
  is consistent with \emph{both} possibilities:
  \begin{itemize}
    \item \textbf{Finitely many twin primes:} The series terminates and
      $B_2$ is a finite sum.
    \item \textbf{Infinitely many twin primes:} The series is infinite
      but still converges, as proved.
  \end{itemize}
\end{remark}

\begin{remark}[Contrast with Euler's theorem]
  Euler proved $\sum_p 1/p = \infty$ precisely because $\pi(x) \sim x/\log x$
  grows fast enough that the sum diverges.  Since $\pi_2(x) = O(x/(\log x)^2)$,
  the twin prime sum has ``one extra logarithm'' in the denominator, which is
  exactly what makes it converge.  The heuristic boundary for convergence of
  $\sum f(p)$ is governed by the size of $\pi_f(x) \sim g(x)$: if
  $g(x)/x \ll (\log x)^{-2}$, convergence is expected.
\end{remark}

\begin{remark}[Progress toward the Twin Prime Conjecture]
  Brun's theorem was the first result to bound $\pi_2(x)$ nontrivially.
  Subsequent progress includes:
  \begin{itemize}
    \item \textbf{Bombieri-Vinogradov} (1965): Large sieve estimates for
      primes in progressions.
    \item \textbf{Chen Jingrun} (1973): There exist infinitely many primes
      $p$ such that $p+2$ is either prime or a product of two primes.
    \item \textbf{Zhang Yitang} (2013): Bounded prime gaps --- there exist
      infinitely many pairs $(p, p')$ with $|p - p'| < 70\,000\,000$.
    \item \textbf{Maynard-Tao} (2014): The gap can be reduced to $246$.
  \end{itemize}
  The Twin Prime Conjecture itself ($\pi_2(x) \to \infty$) remains open.
\end{remark}

% ================================================================
\section{Conclusion}
% ================================================================

We have presented a complete proof of Brun's theorem: the sum of reciprocals
of twin primes converges.  The proof has two components:
\begin{enumerate}
  \item \textbf{Brun's combinatorial sieve} yields the bound
    $\pi_2(x) = O(x(\log\log x)^2/(\log x)^2)$, which is the key
    quantitative input.
  \item \textbf{Partial summation} (Abel's method) converts this bound
    into the convergence of the series.
\end{enumerate}

Brun's theorem is not merely a curiosity: the sieve technique he introduced
has grown into a central branch of number theory.  Modern sieves (Selberg,
large sieve, $\beta$-sieve) are all descendants of Brun's 1919 idea.

\begin{thebibliography}{9}

\bibitem{Brun1919}
V.~Brun.
\newblock La s{\'e}rie $\frac{1}{5}+\frac{1}{7}+\frac{1}{11}+\frac{1}{13}+\frac{1}{17}+\frac{1}{19}+\cdots$
  o\`u les d{\'e}nominateurs sont nombres premiers jumeaux est convergente ou finie.
\newblock \emph{Skrifter udgivne af Videnskabsselskabet i Christiania, Mathematisk-naturvidenskabelig Klasse},
  no.~3 (1919), 9~pp.

\bibitem{HardyWright}
G.~H.~Hardy and E.~M.~Wright.
\newblock \emph{An Introduction to the Theory of Numbers}, 6th~ed.
\newblock Oxford University Press, 2008.

\bibitem{IwaniecKowalski}
H.~Iwaniec and E.~Kowalski.
\newblock \emph{Analytic Number Theory}.
\newblock American Mathematical Society Colloquium Publications, vol.~53,
  Providence, RI, 2004.

\bibitem{CojocMurty}
A.~C.~Cojocaru and M.~R.~Murty.
\newblock \emph{An Introduction to Sieve Methods and Their Applications}.
\newblock Cambridge University Press, 2005.

\bibitem{Nicely1999}
T.~R.~Nicely.
\newblock Enumeration to $1.6 \times 10^{15}$ of the twin primes and Brun's constant.
\newblock \emph{Virginia Journal of Science}, 46(3):195--204, 1995; revised 2010.
\newblock \url{https://faculty.lynchburg.edu/~nicely/twins/twins.html}

\bibitem{OliveiraSilva2014}
T.~Oliveira~e~Silva, S.~Herzog, and S.~Pardi.
\newblock Empirical verification of the even Goldbach conjecture and computation
  of prime gaps up to $4 \times 10^{18}$.
\newblock \emph{Mathematics of Computation}, 83(288):2033--2060, 2014.

\end{thebibliography}

\end{document}
